\selectlanguage{spanish}
Primero que nada, una pequeña disculpa, ya que voy a escribir estos agradecimientos en castellano. También voy a tener que resumir muchísimo para evitar que terminen siendo más largos que la propia tesis.

El doctorado es un camino largo y, por momentos, muy solitario. Sin embargo, a lo largo de estos años tuve la suerte de estar rodeado de gente que estuvo ahí para acompañarme. Sin todes ustedes, nada de esto hubiese sido posible.

No podría empezar sin agradecerle a Diego, quien no solo volvió a demostrar ser un excelente director, sino también una gran persona. Siempre supo escucharme, tanto cuando hablábamos del espectro como cuando necesitaba contarle lo que me estaba pasando, brindándome un apoyo que va muchísimo más allá de su rol académico. Sin él, nada de esto hubiera sido posible.

Quiero agradecerle a mi familia, a Mamá, Papá y Tho, a quienes les dedico esta tesis. Gracias por todo el amor que me dan y por buscar siempre la manera de ayudarme a alcanzar mis objetivos, grandes o chiquitos por igual.

A Mamá, por escucharme siempre, por reírnos cuando encontramos cosas parecidas entre nosotros (que son un montón), por recordar juntos a los abuelos y al pueblo, y por hacer que siempre tenga un pedacito de mi corazón allá, en la precordillera.

A Papá, por enseñarme (aunque todavía debería aprenderlo mejor) a tomarme las cosas con paciencia, pensarlas, desarmarlas, darlas vuelta y tratar de entenderlas. Por sus recomendaciones de caminos, que van mucho más allá de la Ciudad de Buenos Aires, fomentando desde siempre mi amor por la ciencia, incluso congelándose en las salidas de la AAAA cuando era chico y todavía no podía ir solo.

A Tho, que siempre va a ser mi hermano chiquito, por su manera alocada de ver las cosas y por compartir cada día más: libros, juegos, la facultad y, sobre todo, muchísimas risas.

A la Abuela Lucía, que estuvo ahí desde el principio, alentando mi curiosidad. Y también a los abuelos Luis, Bebo y, especialmente, Zunny, a quienes extraño un montón. Al tío Luis, por empujarme a salir de mi zona de confort y ayudarme a estar siempre abierto a descubrir cosas nuevas. Y a Mari, por ser mi hermana del corazón y por haberme hecho padrino de Dante, el mejor ahijado del mundo.

Estar lejos de ustedes fue una de las cosas más difíciles del doctorado. Pero cada día encontré algo que me hizo pensar en ustedes y sentirlos un poquito más cerca.


También quiero agradecer a quienes me acompañaron desde Karlsruhe.

A los investigadores del IAP: a Markus, por las charlas sobre el espectro y nuestro amor compartido por la historia; a David, por los consejos sobre el SD y las charlas al mediodía; a Darko, por sus incontables sugerencias académicas y nuestro amor por el cine; y a Ralph, por su optimismo contagioso y por ir siempre para adelante.

A mis compañerxs de doble doctorado y de oficina, que hicieron que siempre me sintiera en casa. A Stephen, por hacerme reír siempre; a Fiona, por ser tan adorable; a Berenika, por su amor por las pesitas; a Philip, por tantas charlas y horas de trabajo compartidas; a Tobías, por nuestro amor por el terror, la música y las montañas rusas; y a Paul, por ser un gran colega y amigo.

Y muy especialmente a Sari, por ser una persona increíble y por haber logrado que, durante esos seis meses que vivimos juntos, no tuviera solamente una casa, sino un hogar. Cuando quieras, nos volvemos a mudar juntos.

No puedo dejar de agradecer a toda la gente de ITeDA. Cuando empecé a pensar cómo escribir estos agradecimientos, me di cuenta de que tengo anécdotas con prácticamente todes.

A la gente de Ingeniería, por las juntadas a comer, los partidos de fútbol y las charlas de todos los días en el C-101. Al grupo de Auger, por todo el trabajo compartido que, aunque a veces es arduo, nunca viene sin risas mezcladas con rayos cósmicos, especialmente durante los meetings en Malargüe.

Y, especialmente, a mis compañerxs de doble doctorado: Carmi, Marina, Varada, Lourdes, Iván, Ailén y Luca.

A Adri, por estar siempre ahí cuando rompía todo. A Nahui y Juanma, por las salidas compartidas tanto acá como en Karlsruhe. A Fer, por ser un amigo de fierro y por todas las ayudas con las mudanzas. A Joaco, por las aventuras, especialmente en Japón. A José, por las charlas sobre el norte y por el trabajo compartido en el espectro.

A Eze, por ser mi hermano de doctorado. Por todo lo que laburamos juntos, los viajes escuchando C. Tangana y las charlas sobre absolutamente todo.

Y, finalmente, a Belu, por escucharme todos los días, por las lecciones sobre Auger, el doctorado y la vida, y por ser mi amiga desde hace tantos años. Sin vos no hubiese llegado a ITeDA y, sin vos, tampoco hubiese llegado hasta acá.

Quiero agradecer a mis colegas, miembrxs de la colaboración por el constante feedback sobre mi trabajo. Y muy especialmente a lxs miembros del editorial board del paper del espectro Markus, Darko, Piera y Brian. Muchas de sus recomendaciones son ahora parte fundamental de esta tesis.

También soy muy afortunado de tener tantos amigxs y grupos de amigxs que me acompañaron a lo largo de estos años.

A los Bulifandela, que aunque ya no nos veamos tanto como antes, siguen estando siempre ahí. A los HAVs, por tantos momentos compartidos.

A la gente de Amigos de la Astronomía, por traerme siempre preguntas nuevas y hacerme repensar cosas que creía saber. Por dejarme ser, de a ratos, la voz de la asociación y por recordarme constantemente que no alcanza solamente con hacer ciencia: también hay que buscar compartirla con la sociedad. La asociación fue el lugar donde pude cultivar mi amor por la astronomía hasta decidir convertirlo en mi carrera.

A Flo, por nuestro amor compartido por la oscuridad y, sobre todo, por ser la mejor amiga del mundo. Gracias por ser esa persona con la que sé que puedo compartir absolutamente todo.

A Ailén, Clay, Gwen y, especialmente, a Sole, por todas las lecciones aprendidas.

Y a quienes estuvieron presentes a lo largo de todo este tiempo: Tanya, Victoria, Bren, Abi, Diana, Benja, Presti, Meshiael, Melanie, Barbi, Paula, Jannet, Flor, Marina, Alba, Elián, Danelise, Lolo, Nacho, Yas, Ale, y tantas otras personas que hicieron este camino mucho más lindo.

Finalmente, me gustaría agradecerle a mi yo de cuando era chico. Sin ese amor por mirar hacia arriba y perderme entre estrellas, nebulosas y galaxias, probablemente nunca hubiese llegado hasta acá. Me hace muy feliz poder decir que creo que ese chico de diez años estaría orgulloso de la persona en la que me convertí.